% ============================================================================
% regulatory-applicability-EN.tex — SUIT Regulatory Applicability module (EN)
% ============================================================================
%
% STATUS: populated, institution-neutral shared module.
%
% ----------------------------------------------------------------------------
% PURPOSE
% ----------------------------------------------------------------------------
% This is a SHARED, INPUT-ONLY module (no \documentclass): it is \input into a
% host SUIT document that supplies the preamble. The host must provide:
%   * biblatex + \addbibresource{suit-consolidated.bib} (the shared bib symlink)
%   * booktabs, longtable, tabularx, array, multirow
%   * tcolorbox with a "caveatbox" environment (as in the suite documents)
%   * the gls-modes glossary (keys gdpr, nis2, aiact are referenced)
%   * the semantic macros \keyword, \framework, \institution
% No macro is redefined here; if a host lacks the caveatbox environment it can
% supply a one-line \newenvironment{caveatbox}[1][]{\par\noindent}{\par} stub.
%
% The module is jurisdiction-neutral. It frames the size- and
% jurisdiction-INVARIANT functional-obligation spine (F1..F6) as the transferable
% proof, then instantiates that spine per regime from the verified regime
% profiles in docs/_regulatory/*.json, and finally gives an adopter in an
% UNLISTED regime a playbook to derive its own row.
%
% ----------------------------------------------------------------------------
% SOURCING
% ----------------------------------------------------------------------------
% Every per-regime row is sourced from a verified regime profile authored from
% real primary regulation and authoritative secondary sources; every citation
% key used below resolves to a VERIFIED or OK entry in suit-consolidated.bib.
% The mapping is interpretive engineering guidance, NOT legal advice.
% ============================================================================

\section{Regulatory applicability: the invariant obligation spine}
\label{sec:reg-applicability}

\begin{caveatbox}[Status and scope of this module]
\noindent This module records an interpretive engineering mapping
between a small set of functional obligations that every modern
data-and-cyber regulatory regime imposes on a university, and the
\framework{SUIT} control set that an institution can invoke to
discharge them. It is \emph{not} legal advice. The named instruments,
authorities and deadlines are the worked instantiation of each regime
as of mid-2026; an adopting institution must have its mapping reviewed
by qualified counsel and by its own data-protection and
information-security officers before relying on it for a compliance
argument. National transpositions, sectoral overlays and
funder-imposed clauses may add provisions this module does not record.
\end{caveatbox}

The architecture and the playbook of the \framework{SUIT} suite are
not designed against a regulatory checklist; they are designed against
the engineering question of how a university can be
sovereign at the architectural level, and they are then \emph{read}
against the regulations the institution faces. This module performs
that reading at the level of the regime, not the article. Its claim is
deliberately narrow and therefore transferable: across every regime we
examined --- and, we argue, across every regime an adopter will meet
--- the same short list of \keyword{functional obligations} recurs.
The instruments that carry them, the competent authorities that
enforce them and the exact deadlines differ by jurisdiction; the
\emph{functions} do not. That invariant list is the proof an adopter
can carry from one jurisdiction to the next.

\subsection{The functional-obligation spine (F1--F6)}
\label{sec:reg-spine}

We name the spine in role terms so that it survives translation across
regimes. Each obligation is a \emph{function the regulator requires the
institution to perform}, independent of the statute that happens to
require it in a given country.

\begin{description}[leftmargin=2.2em,style=nextline]
  \item[F1 --- Personal-data protection / security of processing.]
    Classify personal data and apply appropriate technical and
    organisational measures to protect its confidentiality, integrity
    and availability. This is the omnibus data-protection duty
    (\gls{gdpr} Art.~32 in the baseline; an equivalent exists in every
    regime examined).
  \item[F2 --- Cybersecurity risk governance.]
    Run a documented, risk-based information-security programme ---
    risk analysis, accountable owner, cyber-hygiene, network and
    information-system security measures --- attaching to the
    \emph{function performed} rather than to headcount.
  \item[F3 --- Incident and breach notification.]
    Detect security and personal-data incidents and notify the
    competent authority (and, where the risk is high, affected
    individuals) within a regime-specific deadline.
  \item[F4 --- Cross-border / international data transfers.]
    Establish a lawful route to move personal data outside the home
    jurisdiction for cloud, identity, research-collaboration and AI
    services.
  \item[F5 --- Electronic identification, trust services and
    AI / technology governance.]
    Provide trustworthy electronic identification, signatures and
    seals, and govern the controlled technologies the institution
    deploys (AI systems; export-controlled or otherwise restricted
    technology) under the applicable risk-based regime.
  \item[F6 --- Accountability, auditable logging and conformity.]
    Retain auditable logs and records and be able to \emph{demonstrate}
    that F1--F5 are in place --- the cross-cutting evidentiary duty
    that makes the other five reviewable.
\end{description}

\paragraph{Why the spine is size-invariant.} The spine is a
\keyword{minimal necessary set}: it is the smallest collection of
functions below which an institution is not compliant in \emph{any} of
the regimes examined, and it does not grow with the size of the
institution. The anchor for the claim is the data-protection regime's
\emph{absence of a de-minimis threshold} for the security,
accountability and breach duties: under the baseline \gls{gdpr}, the
only size-conditioned relief is the Art.~30(5) records-of-processing
derogation below 250 persons, which is itself conditional and does not
touch security, accountability or breach duties~\cite{gdpr}. The
cybersecurity regime attaches the same way: where NIS2 captures an
entity --- by sector-and-size under Art.~2(1), or regardless of size
under Art.~2(2) (e.g. public-administration designation) --- its
Art.~21/23 duties apply in full, with no headcount-scaled
relief~\cite{nis2}. A small
institution therefore performs the same six functions as a large one;
only the \emph{instantiation} --- the number of systems, the size of
the team, the volume of logs --- scales. This is the transferable
core: an adopter that can show it discharges F1--F6 has shown the
regulator-facing minimum, whatever its size and wherever it sits.

\subsection{The obligation \texorpdfstring{$\rightarrow$}{->}
mechanism matrix}
\label{sec:reg-matrix}

The spine becomes operational when each function is mapped to the
\framework{SUIT} control set that contributes to discharging it. The
matrix below is the size- and jurisdiction-invariant core: it is read
identically in every regime, and only the \emph{instrument} column (the
local statute, authority and deadline) is substituted per jurisdiction
in Section~\ref{sec:reg-typology}. The control-set references are
role-based --- they name architectural \emph{capabilities}, not
products --- so the matrix transfers without editing.

\begin{table}[H]
\centering\small
\caption{Obligation~$\rightarrow$~mechanism matrix --- the
invariant core. Each functional obligation maps to the SUIT control
capabilities that contribute to discharging it; only the regulatory
instrument is substituted per jurisdiction.}
\label{tab:reg-matrix}
\renewcommand{\arraystretch}{1.25}
\begin{tabularx}{\textwidth}{%
  >{\hsize=0.42\hsize\linewidth=\hsize\bfseries\raggedright\arraybackslash}X%
  >{\hsize=1.28\hsize\linewidth=\hsize\raggedright\arraybackslash}X%
  >{\hsize=1.30\hsize\linewidth=\hsize\raggedright\arraybackslash}X%
}
\toprule
Functional obligation & SUIT control mechanism (invariant) &
What the regulator can be shown \\
\midrule
F1 Personal-data protection / security of processing
  & Identity, policy, network and cryptographic-services capabilities
    as the technical chain enforcing access control and confidentiality;
    storage-placement constraints; the sovereignty-grid data dimension
  & A documented technical-measures chain proportionate to risk,
    keyed to the data classification \\
F2 Cybersecurity risk governance
  & The ISO/IEC~27005-aligned risk-analysis method; the policy and
    governance capabilities; network segmentation and cyber-hygiene
    baselines; the conformance suite as evidence of effectiveness
  & A risk-based, owner-accountable security programme with current
    risk analysis and tested controls \\
F3 Incident and breach notification
  & Observability capability (logging, monitoring, detection); the
    audit pipeline for timeline reconstruction; the documented
    incident-response procedure with a notification path
  & Detection-to-notification within the regime deadline, with an
    evidenced timeline \\
F4 Cross-border / international data transfers
  & Storage-placement and key-custody controls; the sovereignty-grid
    data dimension recording the disclosure-regime exposure; transfer
    safeguards (encryption, pseudonymisation) as supplementary measures
  & A lawful-transfer route with a documented transfer-risk assessment
    and supplementary technical measures \\
F5 Electronic identification, trust services and
    AI / technology governance
  & Identity and cryptographic-services capabilities (qualified
    signatures/seals); the technology-control and segregation patterns
    for restricted technology; AI-system inventory and human-oversight
    hooks
  & Trustworthy e-identification and a governed inventory of
    controlled technologies with oversight evidence \\
F6 Accountability, auditable logging and conformity
  & The auditable-logging and records capability; the conformance
    suite producing assessment artefacts; the sovereignty grid as the
    design-decision record
  & A standing evidence base demonstrating that F1--F5 are in place
    and reviewable on request \\
\bottomrule
\end{tabularx}
\end{table}

\paragraph{How to use the matrix.} The left column is fixed: it is the
spine of Section~\ref{sec:reg-spine}. The middle column is fixed: it
is the \framework{SUIT} capability set, named in role terms. Only the
binding \emph{instrument} changes per regime --- which statute imposes
F1, which authority enforces F3, what the F3 deadline is. The regime
typology that follows fills exactly that variable part, regime by
regime; the mapping method of Section~\ref{sec:reg-mapping} lets an
adopter in an unlisted regime fill it for itself.

% ============================================================================
\section{Regime typology}
\label{sec:reg-typology}

Regimes differ in \keyword{depth}, not in the spine they impose. A
\keyword{deep} regime (here the EU and the US) carries the spine through
a dense, multi-instrument stack with hard deadlines and conformity
machinery; it earns a full F1--F6 table plus a worked applicability
demonstration. A \keyword{light} regime (here the United Kingdom,
Brazil and South Africa) carries the same spine through a more compact
instrument set; it earns a concise but real F1--F6 table. The five
profiles below span Europe, North America, South America and Africa.
Each is one worked instantiation an adopter substitutes with its own.

% ----------------------------------------------------------------------------
\subsection{European Union / EEA --- deep regime (baseline anchor)}
\label{sec:reg-eu}

The EU/EEA is the \framework{SUIT} baseline anchor: a deep,
directly-applicable stack in which the spine is carried by named
Regulations and a Directive, each with its own competent authority and
operative deadline. The data-protection authority is each Member
State's national DPA, made consistent by the European Data Protection
Board; the cybersecurity authority is the national NIS2 competent
authority and CSIRT, coordinated by ENISA; trust services run through
the national eIDAS supervisory body; AI through national
market-surveillance authorities and the European AI Office; and
transfer adequacy is decided by the European Commission.

\begin{table}[H]
\centering\small
\caption{EU/EEA --- F1--F6 instantiation of the invariant spine.}
\label{tab:reg-eu}
\renewcommand{\arraystretch}{1.25}
\begin{tabularx}{\textwidth}{%
  >{\hsize=0.34\hsize\linewidth=\hsize\bfseries\raggedright\arraybackslash}X%
  >{\hsize=1.46\hsize\linewidth=\hsize\raggedright\arraybackslash}X%
  >{\hsize=1.20\hsize\linewidth=\hsize\raggedright\arraybackslash}X%
}
\toprule
Obligation & Instrument and authority & Deadline / mechanism \\
\midrule
F1 Data protection / security of processing
  & \gls{gdpr} Art.~32~\cite{gdpr}; national DPA, EDPB for consistency
  & Directly applicable since 25 May 2018; standing duty, no size
    threshold \\
F2 Cybersecurity risk governance
  & NIS2 Directive (EU)~2022/2555, with Implementing Regulation
    (EU)~2024/2690 where the institution provides an in-scope digital
    service (otherwise a persuasive technical
    reference)~\cite{nis2,eu-regulation-2024-2690,enisa-risk};
    national NIS2 competent authority / CSIRT, ENISA
  & Transposed by 17 Oct 2024; ten Art.~21(2) measure families
    (MFA, cryptography, access control, supply-chain); function-based
    scope \\
F3 Incident and breach notification
  & \gls{gdpr} Arts.~33--34~\cite{gdpr,edpb-guidelines}; NIS2 Art.~23
    reporting~\cite{nis2}; national DPA and NIS2 authority
  & GDPR: notify DPA without undue delay, $\le$72\,h; NIS2 tiered:
    early warning $\le$24\,h, full notice $\le$72\,h, final report
    $\le$1 month \\
F4 Cross-border transfers
  & \gls{gdpr} Chapter~V: adequacy (Art.~45), SCCs (Art.~46,
    Decision~2021/914), EU--US DPF
    adequacy~\cite{gdpr,eu-sccs2021,eu-dpf-adequacy-2023}; European
    Commission (adequacy), DPAs/EDPB (safeguards)
  & Transfer mechanism: adequacy where it exists, else SCCs + a
    post-\emph{Schrems II} transfer-impact assessment and
    supplementary measures \\
F5 eID, trust and AI governance
  & eIDAS\,2 --- Reg.~(EU)~910/2014 as amended by
    2024/1183~\cite{eidas2-2024}; AI Act
    Reg.~(EU)~2024/1689~\cite{eu-ai-act}; \gls{aiact} via national
    market-surveillance authorities and the European AI Office
  & eIDAS\,2 in force 20 May 2024, EUDI Wallet by end-2026; AI Act
    phased: prohibited practices Feb 2025, GPAI Aug 2025, high-risk
    Aug 2026 (research/R\&D exemptions) \\
F6 Accountability, logging, conformity
  & \gls{gdpr} Arts.~5(2), 24, 30~\cite{gdpr}; NIS2 Art.~21 logging /
    Art.~23 reporting~\cite{nis2,eu-regulation-2024-2690}; national
    single-window platforms~\cite{serima}
  & Standing: records of processing on DPA request; NIS2 logging /
    maturity reporting through the national reporting platform \\
\bottomrule
\end{tabularx}
\end{table}

\paragraph{Worked applicability demonstration (EU).} Take a
mid-sized research-intensive university operating only in one Member
State, with no defence research and a single national campus. Does the
spine apply, and at what depth? \emph{F1} applies in full and
unconditionally: the \gls{gdpr} security duty has no de-minimis
threshold, so the institution must run the technical-measures chain
keyed to its data classification regardless of size~\cite{gdpr}.
\emph{F2} applies because, once the university is in scope --- by
sector-and-size (Art.~2(1)) or by a size-independent route
(Art.~2(2), e.g. public-administration designation) --- the Art.~21
duties apply in full, irrespective of headcount~\cite{nis2}.
\emph{F3} therefore bites at two deadlines simultaneously --- the
\gls{gdpr} 72-hour breach clock to the DPA and the NIS2 tiered clock
(24\,h / 72\,h / 1 month) to the cyber authority --- so the
observability and incident-response capability must satisfy the
\emph{shorter} of the applicable clocks. \emph{F4} engages the moment
the institution uses a non-EEA cloud or collaborates with a
third-country partner: it needs an adequacy route or SCCs plus a
transfer-impact assessment~\cite{eu-sccs2021,eu-dpf-adequacy-2023}.
\emph{F5} engages partially: eIDAS\,2 governs any qualified
signing/sealing of diplomas and contracts~\cite{eidas2-2024}, and the
AI Act engages \emph{only} for deployed (not pure-research) AI
systems, with the high-risk obligations arriving Aug
2026~\cite{eu-ai-act}. \emph{F6} is standing throughout: the auditable
log and records capability is what lets the institution \emph{show} the
DPA and the cyber authority that F1--F5 are in place. The
demonstration shows the method: walk the six functions, mark each as
full / partial / conditional, and record \emph{which} clock or
mechanism governs --- the architecture's contribution is identical in
each case; only the trigger and deadline are read from the instrument
column.

% ----------------------------------------------------------------------------
\subsection{United States --- deep regime (sectoral and fragmented)}
\label{sec:reg-us}

The US carries the same spine through a \emph{sectoral and fragmented}
stack: no omnibus statute, but federal sector statutes (FERPA for
student records, HIPAA for health data, GLBA / FTC Safeguards for
financial-aid data), federal-contracting baselines (NIST
SP~800-171, CMMC, FISMA), export-control regimes (EAR, ITAR) and a
growing state-privacy layer (CCPA/CPRA). Enforcement is distributed
across the Department of Education, HHS Office for Civil Rights, the
FTC, the DoD, Commerce/BIS, State/DDTC and state Attorneys General.

\begin{table}[H]
\centering\small
\caption{United States --- F1--F6 instantiation of the invariant spine.}
\label{tab:reg-us}
\renewcommand{\arraystretch}{1.25}
\begin{tabularx}{\textwidth}{%
  >{\hsize=0.34\hsize\linewidth=\hsize\bfseries\raggedright\arraybackslash}X%
  >{\hsize=1.46\hsize\linewidth=\hsize\raggedright\arraybackslash}X%
  >{\hsize=1.20\hsize\linewidth=\hsize\raggedright\arraybackslash}X%
}
\toprule
Obligation & Instrument and authority & Deadline / mechanism \\
\midrule
F1 Data protection / privacy
  & FERPA (20 U.S.C. \S1232g; 34 C.F.R. Part 99)~\cite{us-ferpa};
    HIPAA Privacy Rule for covered health components; CCPA/CPRA as the
    state-law layer~\cite{us-ccpa-cpra}; ED/SPPO, HHS OCR, state AGs
  & FERPA: no fixed breach deadline, administrative enforcement via
    funding withdrawal; CCPA/CPRA largely exempts universities except
    via vendors / for-profit affiliates \\
F2 Risk governance / written security programme
  & FTC Safeguards Rule (16 C.F.R. Part 314)~\cite{us-ftc-safeguards};
    HIPAA Security Rule risk analysis~\cite{us-hipaa-security}; NIST
    CSF~2.0 (voluntary baseline)~\cite{nist-csf}; FTC, HHS OCR
  & Safeguards full-compliance 9 Jun 2023 (written programme,
    qualified individual, risk assessment); HIPAA risk analysis a
    continuing duty; Title-IV pulls in essentially all universities \\
F3 Incident / breach notification
  & HIPAA Breach Notification Rule (45 C.F.R. \S\S164.400--414)~%
    \cite{us-hipaa-breach}; FTC Safeguards
    notification~\cite{us-ftc-safeguards}; state breach
    statutes~\cite{us-ccpa-cpra}; HHS OCR, FTC, state AGs
  & HIPAA: notify individuals and HHS $\le$60 days (500+ contemporaneous);
    FTC: notify $\le$30 days for events affecting 500+ consumers
    (since 13 May 2024); state timelines layer on top \\
F4 Cross-border transfers
  & Handled functionally through export control rather than a
    transfer-adequacy regime --- see F5; CCPA/CPRA vendor flow-down for
    consumer data~\cite{us-ccpa-cpra}
  & No EU-style adequacy list; the operative cross-border control for
    research data is the deemed-export regime (F5) \\
F5 Trust / AI / technology governance (export control)
  & EAR (15 C.F.R. 730--774, fundamental-research exclusion
    \S734.8)~\cite{us-ear}; ITAR (22 C.F.R. 120--130)~\cite{us-itar};
    Commerce/BIS, State/DDTC
  & Licensing + deemed-export: releasing controlled
    technology/technical data to foreign nationals (even on US soil)
    is an export; operationalised via technology-control plans and
    segregated environments \\
F6 Conformity assessment / certification
  & CMMC Program (32 C.F.R. Part 170) with DFARS
    implementation~\cite{us-cmmc}; NIST SP~800-171 Rev.~3 as the
    assessed control set~\cite{nist-171,nist-171r3,nist-sp800-53};
    DoD (CMMC PMO, DIBCAC), C3PAOs
  & CMMC final rule effective 16 Dec 2024; DFARS clause effective
    10 Nov 2025, phased over ~3 years; Level~2 (CUI) requires a C3PAO
    assessment; artefacts retained six years \\
\bottomrule
\end{tabularx}
\end{table}

\paragraph{Worked applicability demonstration (US).} Take the same
mid-sized research university, now with a Title-IV financial-aid office,
a small health clinic and one DoD-funded research project. The spine
applies, but it is \emph{assembled} from sector statutes rather than
read off one instrument. \emph{F1} is anchored by FERPA over student
records~\cite{us-ferpa}; CCPA/CPRA mostly does \emph{not} reach the
institution directly but reaches its vendors~\cite{us-ccpa-cpra}.
\emph{F2} is triggered the moment the institution participates in
Title-IV: the FTC Safeguards Rule classifies it as a financial
institution and requires the written security programme with a
qualified individual~\cite{us-ftc-safeguards}; the clinic adds the
HIPAA Security Rule risk analysis~\cite{us-hipaa-security}. \emph{F3}
gives the institution \emph{two} breach clocks --- HIPAA's 60-day outer
limit and the FTC's 30-day / 500-consumer trigger --- plus any state
clock; as in the EU case, the architecture must meet the
\emph{shortest} applicable deadline~\cite{us-hipaa-breach}. \emph{F4}
is unusual: there is no adequacy regime, so the cross-border control
that bites is the \emph{deemed-export} concept --- releasing the
DoD-project's controlled technical data to a foreign-national
researcher, even on campus, is an export needing a license or a
technology-control plan~\cite{us-ear,us-itar}. \emph{F5} therefore
\emph{is} the cross-border control here, which is why F4 and F5 fold
together in the US column. \emph{F6} is where the US regime is most
distinctive: CMMC converts the NIST SP~800-171 controls into a
\emph{contractually enforceable, third-party-assessed} certification
that the DoD-funded project cannot be awarded without~\cite{us-cmmc,
nist-171r3}. The demonstration shows the same six functions resolving
to a different instrument geometry --- the spine is invariant, the
assembly is sectoral.

% ----------------------------------------------------------------------------
\subsection{United Kingdom --- light regime (post-Brexit, diverging)}
\label{sec:reg-uk}

The UK carries the spine through retained-and-amended EU-derived
instruments now diverging via domestic reform. The data-protection
authority is the ICO (becoming the Information Commission under the
Data (Use and Access) Act 2025); cyber resilience runs through the NIS
Regulations 2018 with sector competent authorities and the NCSC;
accessibility runs on domestic statute.

\begin{table}[H]
\centering\small
\caption{United Kingdom --- F1--F6 instantiation (light regime).}
\label{tab:reg-uk}
\renewcommand{\arraystretch}{1.25}
\begin{tabularx}{\textwidth}{%
  >{\hsize=0.40\hsize\linewidth=\hsize\bfseries\raggedright\arraybackslash}X%
  >{\hsize=1.42\hsize\linewidth=\hsize\raggedright\arraybackslash}X%
  >{\hsize=1.18\hsize\linewidth=\hsize\raggedright\arraybackslash}X%
}
\toprule
Obligation & Instrument and authority & Deadline / mechanism \\
\midrule
F1 Data protection / security of processing
  & UK GDPR with the Data Protection Act 2018, amended by the Data
    (Use and Access) Act 2025~\cite{uk-gdpr,uk-dpa-2018,uk-duaa-2025};
    ICO / Information Commission
  & Standing security-of-processing duty (Art.~32); DUAA 2025 key
    provisions in force 5 Feb 2026 \\
F2 Cybersecurity risk governance
  & NIS Regulations 2018 (SI 2018/506); forthcoming Cyber Security and
    Resilience Bill~\cite{uk-nis-2018,uk-ncsc,uk-csr-bill}; sector
    competent authorities, ICO (RDSPs), NCSC
  & Risk-management duty; NCSC Cyber Assessment Framework; CSR Bill
    expands scope (e.g. managed service providers), not yet in force \\
F3 Incident / breach notification
  & UK GDPR Art.~33~\cite{uk-gdpr}; NIS Regs reg.~11~\cite{uk-nis-2018};
    ICO; sector competent authority / NCSC
  & GDPR breach to ICO $\le$72\,h; significant NIS incident to the
    competent authority $\le$72\,h \\
F4 Cross-border transfers
  & UK GDPR Chapter~V; ICO IDTA and the UK Addendum to the EU
    SCCs~\cite{uk-idta,uk-gdpr,uk-duaa-2025}; ICO, Secretary of State
    (adequacy)
  & Adequacy where it exists, else IDTA / Addendum + a transfer-risk
    assessment; DUAA 2025 ``not-materially-lower'' data-protection test \\
F5 eID / trust services
  & UK eIDAS (retained Reg.~910/2014 as amended; SI
    2016/696)~\cite{uk-eidas-2016}; ICO as supervisory body
  & Conformity / qualified-status mechanism: ICO grants/revokes
    qualified status with breach reporting and audits \\
F6 Accountability / accessibility conformity
  & Public Sector Bodies Accessibility Regulations 2018 (SI 2018/952),
    WCAG~2.1 AA / EN~301~549~\cite{uk-psbar-2018,uk-jisc-accessibility};
    GDS (monitoring), EHRC (enforcement)
  & Standing conformance + published accessibility statement; GDS
    samples, EHRC may issue unlawful-act notices \\
\bottomrule
\end{tabularx}
\end{table}

% ----------------------------------------------------------------------------
\subsection{Brazil --- light regime (distributed federal stack)}
\label{sec:reg-br}

Brazil carries the spine through a distributed federal stack with no
single omnibus cybersecurity statute: the LGPD (data protection and
breach notification), the national cyber strategy and federal
incident-management decrees, the Inclusion Law (accessibility), the
Access to Information Law (transparency) and the Federal Administrative
Procedure Law (due process). Federal public universities are in scope
as \emph{autarquias / funda\c{c}\~oes p\'ublicas} of the federal public
administration.

\begin{table}[H]
\centering\small
\caption{Brazil --- F1--F6 instantiation (light regime).}
\label{tab:reg-br}
\renewcommand{\arraystretch}{1.25}
\begin{tabularx}{\textwidth}{%
  >{\hsize=0.40\hsize\linewidth=\hsize\bfseries\raggedright\arraybackslash}X%
  >{\hsize=1.42\hsize\linewidth=\hsize\raggedright\arraybackslash}X%
  >{\hsize=1.18\hsize\linewidth=\hsize\raggedright\arraybackslash}X%
}
\toprule
Obligation & Instrument and authority & Deadline / mechanism \\
\midrule
F1 Data protection / security of processing
  & LGPD (Lei n\textsuperscript{o}~13.709/2018), esp. Art.~46~%
    \cite{br-lgpd-13709}; ANPD
  & Continuous duty: technical and administrative security measures;
    ANPD may set minimum standards; applies to public-power controllers \\
F2 Cybersecurity governance / risk management
  & E-Ciber national cyber strategy (Decreto
    n\textsuperscript{o}~10.222/2020)~\cite{br-eciber-d10222};
    Marco Civil da Internet baseline~\cite{br-marco-civil-12965};
    GSI/PR
  & Federal bodies (incl. universities) implement strategic actions
    under GSI coordination; strategy-cycle instrument plus GSI
    normative instructions \\
F3 Incident / breach notification
  & LGPD Art.~48 operationalised by Resolu\c{c}\~ao CD/ANPD
    n\textsuperscript{o}~15/2024~\cite{br-anpd-res15-2024}; federal
    cyber track via Decreto n\textsuperscript{o}~10.748/2021
    (CTIR.Gov)~\cite{br-redeinc-d10748}; ANPD, GSI/PR--CTIR.Gov
  & Personal-data breaches: notify ANPD and subjects $\le$3 business
    days (supplementable $\le$20 business days); 5-year retention;
    federal cyber incidents via CTIR.Gov \\
F4 Cross-border transfers
  & LGPD international-transfer regime (Chapter~V of Lei
    13.709/2018)~\cite{br-lgpd-13709}; ANPD
  & Transfer on an adequacy / safeguards / consent / contractual
    basis; ANPD supervises (no EU-style adequacy list) \\
F5 Accessibility / inclusion (trust-equivalent duty)
  & Lei Brasileira de Inclus\~ao (Lei n\textsuperscript{o}~13.146/2015),
    Arts.~63\,ff., operationalised by e-MAG~\cite{br-lbi-13146}; MEC,
    Governo Digital
  & Standing accessibility duty assessed against e-MAG; conditions MEC
    course authorisation/recognition \\
F6 Accountability / transparency / due process
  & Access to Information Law (Lei n\textsuperscript{o}~12.527/2011)~%
    \cite{br-lai-12527}; Federal Administrative Procedure Law (Lei
    n\textsuperscript{o}~9.784/1999)~\cite{br-procadm-9784}; CGU,
    each federal body
  & Active + passive disclosure with statutory response deadlines;
    baseline procedural guarantees governing decisions arising from
    F1--F5 \\
\bottomrule
\end{tabularx}
\end{table}

% ----------------------------------------------------------------------------
\subsection{South Africa --- light regime (single national jurisdiction)}
\label{sec:reg-za}

South Africa carries the spine through a single national jurisdiction
with no member-state transposition layer and no standalone unified
Cybersecurity Act. POPIA supplies data protection, breach notification
and cross-border transfer; the Cybercrimes Act supplies incident /
offence reporting; the ECT Act supplies trust services; and
accessibility is principle-based under the Equality Act. The African
Union's Malabo Convention is regional \emph{context} only --- South
Africa has signed but not ratified it, so no transfer obligation flows
from it to ZA institutions~\cite{za_malabo}.

\begin{table}[H]
\centering\small
\caption{South Africa --- F1--F6 instantiation (light regime).}
\label{tab:reg-za}
\renewcommand{\arraystretch}{1.25}
\begin{tabularx}{\textwidth}{%
  >{\hsize=0.40\hsize\linewidth=\hsize\bfseries\raggedright\arraybackslash}X%
  >{\hsize=1.42\hsize\linewidth=\hsize\raggedright\arraybackslash}X%
  >{\hsize=1.18\hsize\linewidth=\hsize\raggedright\arraybackslash}X%
}
\toprule
Obligation & Instrument and authority & Deadline / mechanism \\
\midrule
F1 Data protection / security of processing
  & POPIA (Act 4 of 2013), Condition~7,
    s.19~\cite{za_popia,za_inforegulator}; Information Regulator
  & Continuous risk-based duty: appropriate, reasonable technical and
    organisational measures; no threshold / size carve-out \\
F2 Cybersecurity risk governance
  & POPIA Condition~7 security safeguards read with the Cybercrimes
    Act baseline~\cite{za_popia,za_cybercrimes}; Information Regulator,
    SAPS
  & Risk-based security duty; no standalone unified Cybersecurity Act,
    so the security limb is carried by POPIA \\
F3 Incident / breach notification
  & POPIA s.22~\cite{za_popia}; Cybercrimes Act s.54
    (offence reporting)~\cite{za_cybercrimes}; Information Regulator,
    SAPS
  & POPIA: notify Regulator and subjects as soon as reasonably
    possible after discovery (no hour-count); Cybercrimes Act: report
    qualifying offence $\le$72\,h where the university is an ECS provider \\
F4 Cross-border transfers
  & POPIA s.72~\cite{za_popia,za_inforegulator}; Information Regulator
    (oversight), responsible party self-assesses adequacy
  & Transfer on adequate-protection / BCR / binding-agreement /
    consent / contractual-necessity basis; adequacy self-assessed and
    documented per transfer (no adequacy list) \\
F5 eID / trust services
  & ECT Act (Act 25 of 2002), ss.13 and
    37--40~\cite{za_ect_act}; South African Accreditation Authority
  & Legal recognition of data messages and signatures; advanced
    electronic signature requires an accredited provider's certificate
    (rebuttable presumption of validity) \\
F6 Accountability / accessibility (non-discrimination)
  & PEPUDA / Equality Act (Act 4 of 2000), s.9~\cite{za_pepuda};
    Equality Courts, SAHRC
  & Standing reasonable-accommodation duty; no dedicated
    web-accessibility statute, WCAG as de-facto benchmark \\
\bottomrule
\end{tabularx}
\end{table}

% ============================================================================
\section{Mapping method: deriving a row for an unlisted regime}
\label{sec:reg-mapping}

The five profiles above are worked instantiations, not an exhaustive
list. An adopter in any other jurisdiction derives its own row by the
same procedure: the spine of Section~\ref{sec:reg-spine} is fixed, the
control set of Table~\ref{tab:reg-matrix} is fixed, and only the
\emph{instrument} column has to be filled. The playbook below fills it.

\begin{description}[leftmargin=2.2em,style=nextline]
  \item[Step 1 --- Identify your data-protection authority and its
    security duty (F1).]
    Find the competent authority for personal data (a national DPA, a
    sectoral regulator, or --- as in the US --- several). Locate its
    \emph{security-of-processing} duty and confirm whether a
    size/de-minimis threshold exists. In every regime examined, the
    security duty has \emph{no} size exemption; treat its absence as
    the default and verify it explicitly. Map the duty to the F1 row of
    Table~\ref{tab:reg-matrix}.
  \item[Step 2 --- Identify your lawful-basis and risk-governance
    regime (F2).]
    Determine the lawful basis on which the institution processes
    personal data (public task, legitimate interest, consent) and the
    statute or framework that requires a documented, risk-based
    security programme. Where no dedicated cyber statute exists (as in
    Brazil or South Africa), the risk-governance duty is usually
    carried by the security limb of the data-protection law; record
    which instrument carries it.
  \item[Step 3 --- Identify your breach / incident deadline (F3).]
    Find every applicable notification clock --- the data-protection
    breach clock and any cyber-incident clock --- and record the
    \emph{shortest} one, because the observability and
    incident-response capability must satisfy the tightest deadline
    that can bite. Note both the recipient (DPA, CSIRT, police) and the
    trigger (risk-of-harm, affected-count, offence type).
  \item[Step 4 --- Identify your essential-entity / cyber-duty status
    (F2/F3 scope).]
    Determine whether the institution falls into an essential- or
    important-entity (or equivalent) category, and on what basis ---
    function-performed (as under NIS2) or sector listing. This sets the
    depth at which F2 and F3 apply and whether tiered (24\,h / 72\,h /
    final) reporting is in scope.
  \item[Step 5 --- Identify your transfer regime (F4).]
    Establish how personal data may lawfully leave the jurisdiction:
    an adequacy-list mechanism (EU/UK style), self-assessed adequacy
    plus safeguards (Brazil / South Africa style), or a functional
    substitute such as export control (US style). Record the documented
    artefact the regime expects (transfer-impact assessment, binding
    agreement, technology-control plan).
  \item[Step 6 --- Identify your sector overlays (F5/F6).]
    Add the sector-specific layers the institution triggers: a
    trust-services / eID regime for signed diplomas and contracts; an
    AI-governance regime for deployed AI; export / controlled-technology
    rules for restricted research; health, financial-aid or
    accessibility statutes for the relevant units. Each overlay maps to
    F5 (governed technology) or F6 (conformity / accountability).
  \item[Step 7 --- Bind each finding to the control set.]
    For each of Steps~1--6, write the local instrument, authority and
    deadline into the instrument column against the corresponding F
    row, leaving the SUIT control-mechanism column of
    Table~\ref{tab:reg-matrix} unchanged. The result is the adopter's
    own regime profile in the same shape as the five worked profiles
    above --- and the same architectural evidence base discharges it.
\end{description}

\paragraph{Closing note on transferability.} The procedure terminates
because the spine is finite and size-invariant: there are six functions,
each maps to a fixed control capability, and the only per-jurisdiction
work is reading the local instrument that carries each function. An
adopter that completes Steps~1--7 has reproduced, for its own regime,
the proof the five worked profiles already carry --- that discharging
F1--F6 is the regulator-facing minimum, and that the
\framework{SUIT} control set is the same in every jurisdiction. The
mapping is interpretive engineering guidance and must be confirmed
against the adopter's own legal framework and counsel before being
relied upon.

\endinput
